\chapter{INTRODUCTION}
\thispagestyle{empty}

\section{Portée du projet}

Ce projet étudie la recommandation d'articles Wikipedia dans le domaine de
l'informatique. L'objectif à long terme n'est pas seulement de retrouver des
articles liés, mais aussi de les organiser en parcours d'apprentissage utiles
pour un lecteur. Le projet combine donc trois familles de composants :
\begin{itemize}
    \item des plongements sémantiques de texte obtenus à partir du contenu des
    articles,
    \item des plongements structurels de graphe appris à partir de la structure
    des hyperliens de Wikipedia,
    \item des modèles de référence fondés sur les réseaux de neurones sur
    graphes et une future couche de présentation orientée parcours
    d'apprentissage.
\end{itemize}

Le socle expérimental du travail reste le jeu de données Custom-WikiCS, qui
étend le benchmark Wiki-CS~\cite{mernyei2020wikics} avec le texte des
articles, des caches de plongements et des artefacts d'évaluation adaptés à
des expériences orientées recommandation. Dans la branche actuelle, le projet
utilise les plongements textuels BGE-M3~\cite{chen2025bgem3}, des plongements
structurels de type Node2Vec, une base de référence Graph Convolutional Network
(GCN)~\cite{kipf2017semi} et une base de référence GraphSAGE~\cite{hamilton2017inductive}.
Les Graph Attention Networks (GAT)~\cite{velickovic2018graph} restent présents
dans l'orientation conceptuelle du projet, mais les résultats les plus stables
à ce stade proviennent de la branche d'évaluation révisée centrée sur une
recherche hybride sémantique-structurelle.

\section{État du projet au début du troisième rapport intermédiaire}

Au début de ce troisième rapport intermédiaire, le projet a largement dépassé
le stade de la définition du problème et de la construction du jeu de données.
Les jalons majeurs suivants sont déjà atteints :
\begin{itemize}
    \item construction d'un graphe d'articles Wikipedia nettoyé comportant
    11\,701 nœuds et 291\,039 arêtes orientées,
    \item génération des plongements textuels BGE-M3 et de plusieurs branches
    structurelles fondées sur Node2Vec,
    \item implémentation des bases de référence structurelles GCN et
    GraphSAGE,
    \item développement d'un prototype local d'application web pour la
    recherche interactive d'articles,
    \item correction des problèmes d'évaluation antérieurs et création d'une
    branche structurelle non orientée révisée pour une comparaison plus juste.
\end{itemize}

L'objectif principal de ce rapport est de présenter les premiers résultats
consolidés de cette branche révisée, d'analyser ce qui a changé par rapport
aux expériences antérieures, de documenter les problèmes rencontrés au cours
du processus et de formaliser les métriques de performance ainsi que le
protocole d'évaluation actuellement utilisés pour la vérification du projet.
